% ============================================================
% FRONT MATTER — Lời nói đầu, Hướng dẫn sử dụng
% ============================================================

\frontmatter

% --- Mục lục ---
\tableofcontents
\clearpage

% --- Lời nói đầu ---
\chapter*{Lời nói đầu}
\addcontentsline{toc}{chapter}{Lời nói đầu}

Nhằm mục đích hỗ trợ các học viên trong việc chuẩn bị cho kỳ thi IELTS, \textbf{The Forum Center} đã ra mắt dự án \textit{IELTS Speaking Forecast Intensive}.

Tài liệu này tập trung vào việc tổng hợp và giải chi tiết các câu hỏi dự đoán (Forecast) bám sát xu hướng đề thi thực tế mới nhất, bao gồm:
\begin{itemize}[leftmargin=2em, itemsep=4pt]
  \item \textbf{Part 1}: Các chủ đề xoay quanh cuộc sống thường ngày.
  \item \textbf{Part 2}: Các dạng bài mô tả trọng tâm (Cue Cards) rèn luyện phát triển ý tưởng trong 2 phút.
  \item \textbf{Part 3}: Các câu hỏi nghị luận xã hội đòi hỏi tư duy sâu sắc và cách trả lời logic.
\end{itemize}

Mỗi chủ đề đều được thiết kế với dàn ý rõ ràng, cung cấp ý tưởng mở rộng, từ vựng ăn điểm và bài mẫu mang tính ứng dụng cao.

\vspace{10pt}
\begin{center}
  \textit{\color{ForumDarkBlue}Chúc các bạn ôn tập hiệu quả và chinh phục thành công mục tiêu IELTS!}
\end{center}

\vfill
\begin{flushright}
  \textbf{Phòng Đào Tạo — The Forum Center}
\end{flushright}

\clearpage

% --- Hướng dẫn sử dụng ---
\chapter*{Hướng dẫn sử dụng}
\addcontentsline{toc}{chapter}{Hướng dẫn sử dụng}

\begin{prompt}[title={Cách tiếp cận bài học}]
  \begin{enumerate}[leftmargin=2em, itemsep=4pt]
    \item \textbf{Lên ý tưởng}: Tự phác thảo ý chính dựa trên câu hỏi trước khi tham khảo câu trả lời mẫu.
    \item \textbf{Đọc hiểu \& Phân tích}: Nghiên cứu các cụm từ (collocations) và cách tổ chức ý trong bài giải.
    \item \textbf{Thực hành phát âm}: Đọc thành tiếng các câu trả lời, chú ý đến các ngắt nghỉ, nhấn nhá đã được gợi ý.
    \item \textbf{Cá nhân hoá}: Áp dụng từ vựng để tạo câu trả lời của riêng mình.
  \end{enumerate}
\end{prompt}
